% Generated by mceg equations; edit the contract, then regenerate.
% Requires amsmath. Authored math is retained; TeX is not executed here.
% Definitions and hashes do not establish scientific validity.

% equation-sha256: 5f1b83db451a1beddd5bebac5695edc4e41446ade8bd628210b6d0fbbab65238
\paragraph{Equation wind\_momentum (revision 1)}
Pressure, steady spherical continuity and prescribed-temperature momentum balance.
\begin{equation}
\begin{aligned}
a^2=\frac{k_{\rm B}T(r)}{\mu m_{\rm p}},\quad P=\rho a^2,\quad 4\pi r^2\rho v=\dot M,\\
\kappa=\frac{3GM_\star}{a_{\rm orb}^3},\quad g_{\rm eff}=\frac{GM_{\rm p}}{r^2}-\kappa r,\\
\left(v-\frac{a^2}{v}\right)\frac{dv}{dr}=\frac{2a^2}{r}-\frac{da^2}{dr}-g_{\rm eff}
\end{aligned}
\label{eq:escape.parker_5f_wind:wind_5f_momentum:1}
\end{equation}
\noindent\textbf{Scope:} Constant mu; optional linear tide along the star-planet axis, combined with spherical continuity.

\begin{description}
\item[$r$] Radius (cm; ID: radius).
\item[$T$] Prescribed temperature (K; ID: temperature).
\item[$\mu$] Mean particle mass in proton-mass units (1; ID: mu).
\item[$M_{\rm p}$] Planet mass (g; ID: mass).
\item[$\kappa$] Supplied tidal coefficient (s\textasciicircum{}-2; ID: tide).
\item[$v$] Internal radial velocity (cm s\textasciicircum{}-1; ID: velocity).
\item[$a$] Local isothermal sound speed (cm s\textasciicircum{}-1; ID: sound\_speed).
\item[$\rho$] Mass density; normalization not determined by the velocity solver (g cm\textasciicircum{}-3; ID: density).
\end{description}

\noindent\textbf{Assumptions:}
\begin{itemize}
\item The caller supplies kappa; the wind API does not infer stellar mass or orbital separation.
\item The stellar-mass-dominated, r much smaller than orbital separation expansion is not full Roche geometry.
\item Temperature is supplied and mu is fixed. No energy or ionization closure is solved.
\end{itemize}

% equation-sha256: 9d848f938218e7c751318220f2305fc9c64d21b847ff2a201a66de2b90d3dcca
\paragraph{Equation isothermal\_limit (revision 1)}
No-tide isothermal Parker branch and its regular derivative.
\begin{equation}
\begin{aligned}
r_{\rm s}=\frac{GM_{\rm p}\mu m_{\rm p}}{2k_{\rm B}T},\quad t=\ln(r/r_{\rm s}),\quad z=\ln(v/a),\\
e^{2z}-2z-1=4(t+e^{-t}-1),\quad \operatorname{sign}z=\operatorname{sign}t,\\
\frac{dz}{dt}=\frac{2(1-e^{-t})}{e^{2z}-1},\qquad \left.\frac{dz}{dt}\right|_{t=0}=1
\end{aligned}
\label{eq:escape.parker_5f_wind:isothermal_5f_limit:1}
\end{equation}
\noindent\textbf{Scope:} The analytic-limit implementation solves the scalar implicit relation by bisection; it does not call Lambert W.

\begin{description}
\item[$r$] Radius (cm; ID: radius).
\item[$r_{\rm s}$] Isothermal sonic radius (cm; ID: sonic\_radius).
\item[$T$] Constant temperature (K; ID: temperature).
\item[$M_{\rm p}$] Planet mass (g; ID: mass).
\item[$\mu$] Mean particle mass in proton-mass units (1; ID: mu).
\item[$t$] Log radius divided by sonic radius (1; ID: coordinate).
\item[$z$] Log Mach number (1; ID: log\_mach).
\end{description}

\noindent\textbf{Assumptions:}
\begin{itemize}
\item Float64 and positive inputs; the accelerating wind branch is selected.
\item A fifth-order sonic series for |t|<1e-3 avoids cancellation; custom JVP uses its derivative.
\end{itemize}

% equation-sha256: b9620ae024389487663d0f4c42b673fa718090a82bb0692eb89577063e07a8f3
\paragraph{Equation dimensionless\_wind (revision 1)}
Dimensionless residual and continuation from the isothermal no-tide wind.
\begin{equation}
\begin{aligned}
x=\ln(r/R),\quad L=\ln(r_{\rm out}/R),\quad a_{\rm ref}^2=\frac{k_{\rm B}T(r_{\rm out})}{\mu m_{\rm p}},\\
s_\lambda=1+\lambda\left[\frac{T(Re^x)}{T(r_{\rm out})}-1\right],\quad \beta_\lambda=\frac{\partial\ln s_\lambda}{\partial x},\\
A=\frac{GM_{\rm p}}{Ra_{\rm ref}^2},\quad B=\frac{\kappa R^2}{a_{\rm ref}^2},\quad q=\ln\frac{v}{a_{\rm ref}\sqrt{s_\lambda}},\\
N_\lambda=2-\beta_\lambda-\frac{Ae^{-x}}{s_\lambda}+\frac{\lambda Be^{2x}}{s_\lambda},\\
(e^{2q}-1)\left(q_x+\frac{\beta_\lambda}{2}\right)-N_\lambda=0
\end{aligned}
\label{eq:escape.parker_5f_wind:dimensionless_5f_wind:1}
\end{equation}
\noindent\textbf{Scope:} The target model is lambda=1. Intermediate lambda values are numerical initialization only.

\begin{description}
\item[$r$] Radius (cm; ID: radius).
\item[$R$] Inner domain radius (cm; ID: base).
\item[$r_{\rm out}$] Outer domain radius (cm; ID: outer).
\item[$x$] Log radius relative to the inner boundary (1; ID: coordinate).
\item[$\lambda$] Continuation fraction (1; ID: continuation).
\item[$s_\lambda$] Squared sound speed divided by its reference value (1; ID: sound\_ratio).
\item[$q$] Log Mach number (1; ID: log\_mach).
\item[$N_\lambda$] Dimensionless momentum numerator (1; ID: numerator).
\end{description}

\noindent\textbf{Assumptions:}
\begin{itemize}
\item Eight continuation steps; default 129 total collocation nodes.
\item Reference temperature is evaluated at the outer boundary, not at the sonic point.
\end{itemize}

% equation-sha256: 722d8907a85c6c31aa007247f25cd41aa148bbbda06f04bb3dfdb7fe39c7cd4f
\paragraph{Equation sonic\_regularity (revision 1)}
Regular transonic critical point and slope.
\begin{equation}
\begin{aligned}
N_\lambda(x_{\rm c})=0,\quad q(x_{\rm c})=0,\quad N_{\lambda,x}(x_{\rm c})>0,\\
q_{x,{\rm c}}=\frac{-\beta_{\lambda,{\rm c}}+\sqrt{\beta_{\lambda,{\rm c}}^2+8N_{\lambda,x}(x_{\rm c})}}{4},\\
\left.\frac{d\ln v}{d\ln r}\right|_{\rm c}=q_{x,{\rm c}}+\frac{\beta_{\lambda,{\rm c}}}{2}
\end{aligned}
\label{eq:escape.parker_5f_wind:sonic_5f_regularity:1}
\end{equation}
\noindent\textbf{Scope:} One negative-to-positive numerator root in the supplied radial interval.

\begin{description}
\item[$x_{\rm c}$] Log critical radius (1; ID: critical).
\item[$N_\lambda$] Dimensionless numerator (1; ID: numerator).
\item[$\beta_\lambda$] Log sound-speed-squared slope (1; ID: beta).
\item[$q$] Log Mach number (1; ID: log\_mach).
\end{description}

\noindent\textbf{Assumptions:}
\begin{itemize}
\item The implementation also requires subsonic inner and supersonic outer branches.
\item The no-tide isothermal initializer critical point must lie inside the domain.
\end{itemize}

% equation-sha256: 3c88e12900ef54fcc8caeed05c93ea579519ec2b18a330d6ed45a7cfc1098ff6
\paragraph{Equation collocation (revision 1)}
Two Chebyshev-Lobatto pieces joined at the critical point.
\begin{equation}
\begin{aligned}
\xi_i=\frac{1-\cos(\pi i/(m-1))}{2},\quad i=0,\ldots,m-1,\\
x_i^-=x_{\rm c}\xi_i,\qquad x_i^+=x_{\rm c}+(L-x_{\rm c})\xi_i,\\
\mathcal R_i=(e^{2q_i}-1)\left[(D_xq)_i+\frac{\beta_i}{2}\right]-N_i=0,\\
\mathcal R_{\rm c}\leftarrow q_{\rm c},\quad
\mathcal R_0\leftarrow(D_x^-q)_{\rm c}-q_{x,{\rm c}},\quad
\mathcal R_{n-1}\leftarrow(D_x^+q)_{\rm c}-q_{x,{\rm c}}
\end{aligned}
\label{eq:escape.parker_5f_wind:collocation:1}
\end{equation}
\noindent\textbf{Scope:} The repeated critical node is stored once: total n=2m-1. The displayed first/last replacements are residual rows, not imposed endpoint velocities.

\begin{description}
\item[$n$] Odd total node count; at least 17 (1; ID: grid\_size).
\item[$x$] Log radial coordinate (1; ID: coordinate).
\item[$x_{\rm c}$] Log critical radius (1; ID: critical).
\item[$D_x$] Barycentric derivative matrix rescaled to each log-radius interval (1; ID: derivative).
\item[$q$] Collocated log Mach number (1; ID: log\_mach).
\item[$\mathcal R$] Discrete momentum and regularity residual (1; ID: residual).
\end{description}

\noindent\textbf{Assumptions:}
\begin{itemize}
\item A damped Newton solve follows continuation; accepted residual infinity norm must be below 1e-8.
\item This is a discretized ODE, not an exact continuum solution.
\end{itemize}

% equation-sha256: d9d27bdf5b0169c836c035a3c92462862d78924cac116031f465ebde2a82fce6
\paragraph{Equation implicit\_derivative (revision 1)}
Implicit derivatives of the critical point and the converged discrete wind.
\begin{equation}
\begin{aligned}
\frac{dx_{\rm c}}{d\theta}=-\frac{N_\theta}{N_x},\qquad
\frac{dq}{d\theta}=-\mathcal R_q^{-1}\left(\mathcal R_\theta+\mathcal R_{x_{\rm c}}\frac{dx_{\rm c}}{d\theta}\right)
\end{aligned}
\label{eq:escape.parker_5f_wind:implicit_5f_derivative:1}
\end{equation}
\noindent\textbf{Scope:} At lambda=1; partial derivatives include prescribed-temperature parameters, mass, mu, tide and domain bounds.

\begin{description}
\item[$\theta$] One physical input parameter; units depend on parameter (parameter dependent; ID: parameters).
\item[$x_{\rm c}$] Log critical radius (1; ID: critical).
\item[$q$] Vector of collocated log Mach numbers (1; ID: log\_mach).
\item[$\mathcal R$] Discrete residual including both moving grids (1; ID: residual).
\end{description}

\noindent\textbf{Assumptions:}
\begin{itemize}
\item jax.lax.custom\_root differentiates the residual equations, not bisection or Newton iterations.
\item The scalar critical slope and the discrete residual Jacobian must be nonsingular.
\item The output interpolation and sound-speed factor are differentiated as well.
\end{itemize}

% equation-sha256: 897616fa9ca51dd9bc6361dede83dc37bf52202d037e3b175b51486a3f95db10
\paragraph{Equation velocity\_output (revision 1)}
Public output converts cgs sound speed and velocity to km/s.
\begin{equation}
\begin{aligned}
v_{\rm out}(r)=10^{-5}\sqrt{\frac{k_{\rm B}T(r)}{\mu m_{\rm p}}}\,\exp q(\ln(r/R))\quad[\mathrm{km\,s^{-1}}]
\end{aligned}
\label{eq:escape.parker_5f_wind:velocity_5f_output:1}
\end{equation}
\noindent\textbf{Scope:} For the analytic isothermal function, use z(ln(r/r\_s)) in place of q(ln(r/R)), with constant T.

\begin{description}
\item[$r$] Requested radius (cm; ID: radius).
\item[$T$] Prescribed or constant temperature (K; ID: temperature).
\item[$\mu$] Mean particle mass in proton-mass units (1; ID: mu).
\item[$v_{\rm out}$] Returned radial velocity (km s\textasciicircum{}-1; ID: velocity\_kms).
\item[$q$] Interpolated log Mach number (1; ID: log\_mach).
\end{description}

\noindent\textbf{Assumptions:}
\begin{itemize}
\item Nonisothermal values use piecewise Chebyshev interpolation inside the domain.
\item Invalid or unconverged nonisothermal solutions return NaN; diagnostics are available with return\_info=True.
\end{itemize}
